% =====================================================================
% 9. DISCUSSION                     target: 2 - 2.5 pages
%
% Section 8 says WHAT happened. This section says WHAT IT MEANS.
% This is where personal analysis and scientific judgement are shown,
% so it weighs heavily in the evaluation. Use "I" throughout.
% =====================================================================

\section{Analysis and Discussion}
\label{sec:discussion}

\begin{guide}
Four subsections: what I found, what did not work, why the numbers must be read
carefully, and the limits of the study. Keep it honest -- an internship report that
admits a negative result reads as more competent, not less.
\end{guide}

% ---------------------------------------------------------------------
\subsection{What worked}
\label{sec:disc:worked}

\begin{guide}
Three or four short paragraphs, one finding each. Structure every one the same way:
the observation $\rightarrow$ the evidence (point at a table) $\rightarrow$ my
explanation of \emph{why} it happens.
\textbf{Questions to answer:} Which of my design goals (\Cref{sec:arch:goals}) were
actually reached? Which head repaired the weakness it was designed for? Does the
mechanism I claimed in \Cref{sec:pooling} show up in the numbers, or did the head
help for a different reason?
\end{guide}

\todo{Finding 1 -- the size / accuracy trade-off.}

\todo{Finding 2 -- explanation quality.}

\todo{Finding 3 -- which encoder $\times$ head pairing wins and why.}

% ---------------------------------------------------------------------
\subsection{What did not work}
\label{sec:disc:negative}

\begin{guide}
A real section, not an apology. Every honest study has these, and reporting them is
part of the scientific method.
\textbf{Questions to answer:} Which of my ideas looked good on paper and did not pay
off? Do I know why, or is it still open? What would I test next to find out?
\end{guide}

\todo{Negative result 1 -- e.g. a head that generalises the baseline but does not beat
it in practice. Say what the theory promised (\Cref{prop:ad} or similar) and why the
promise did not materialise: optimisation difficulty, not enough data, the weakness it
fixes was not the binding one.}

\todo{Negative result 2.}

\todo{One sentence on what these negative results teach: extra capacity is only useful
if training can find it.}

% ---------------------------------------------------------------------
\subsection{How the numbers should be read}
\label{sec:disc:reading}

\begin{guide}
This subsection protects my own credibility. Two points must be made clearly:
\begin{itemize}\itemsep2pt
  \item AOPCR is not normalised, so its scale follows the scale of the logits. My
        values must be compared with each other, never with the published column.
        The proof that this is a measurement-scale effect and not a model effect is
        that my \emph{re-trained} InceptionTime shows the same shift.
  \item Any difference smaller than the seed-to-seed standard deviation is not a
        result. Say what that spread is, and refuse to claim wins inside it.
\end{itemize}
\end{guide}

\todo{Paragraph 1 -- the AOPCR scale issue, with the re-trained baseline as evidence.}

\todo{Paragraph 2 -- seed variance and what counts as a real difference.}

\todo{Paragraph 3 -- any difference between my training budget and the published one
(number of repeats, epochs), and how that shifts the comparison.}

% ---------------------------------------------------------------------
\subsection{Limitations}
\label{sec:disc:limitations}

\begin{guide}
Bullet list, three to five items, one line each. Be specific: ``synthetic data'' is
weak, ``\webtraffic{} is synthetic, so the per-time-step ground truth is exact but the
signal is easier than real traffic'' is strong.
\end{guide}

\begin{itemize}\itemsep2pt
  \item \todo{\webtraffic{} is synthetic; it is the only dataset with per-time-step
        ground truth, so explanation quality is measured on one dataset only.}
  \item \todo{Number of seeds and what that means for the size of a claim.}
  \item \todo{Only univariate series were tested, although the architecture supports
        multivariate input.}
  \item \todo{Some variants were screened on \webtraffic{} alone before the full
        \ucr{} sweep, which biases the selection slightly.}
  \item \todo{No human study: a high AOPCR means the explanation is faithful to the
        model, not that a person finds it useful.}
\end{itemize}
